\scanpage{18}\mspage{17}
% @ opening
Une \uline{contrée} est définie par

% @ 1°
1°) Un ensemble de ``figures élémentaires'', $\mathcal{F}_{el}$

% @ 2°
2°) Une relation d'ordre $\ll$ \deleted{$\unsure{\beta|b}$} (raffinement) parmi les figures élémentaires

% @ 3°
3°) Une relation \unsure{reflexive|réflexive} \unsure{symetrique|symétrique} $\mathcal{R}$ (\deleted{\unsure{admissibles|admissibilité|constructibles}}\inserted{``compatibilité''}) parmi figures élémentaires

% @ NB, oblique from left across 1°–4°
% slanted
\margin{\uline{NB} Il revient au $=$ de se donner 4°), ou de se donner la notion de \uline{figure} [une \deleted{partie} configuration est une figure si, en \unsure{même} \unsure{tps} que \unsure{chaque|toute} figure él. qu'elle \unsure{contient}, elle \unsure{contient} \unsure{aussi} \unsure{toute} \unsure{figure} $\leq$ une qu'elle \unsure{englobe|engendre|contient}] et la propriété pour une figure d'être \unsure{finie}.}

% @ 4°
4°) Un ens de \unsure{parties|partie} de $\mathcal{F}_{el}$, appelées les \unsure{---} finies -- dans le cas d'une ``contrée \uuline{finie}'' ce sont les parties \uline{finies} qui sont des ``configurations'' \unsure{ou} i.e. \unsure{formées|familles|famés} de figures él. deux : deux compatibles
% @ small note at left of 4°
% slanted
\margin{\uline{contrée} \uline{finie} ou plutôt, ``\uline{compacte}'' : celle où toute figure est finie}
Satisfaisant : des axiomes convenables, à dégager. Il y a déjà celui-ci

% @ condition 1
\uline{\unsure{Contr|Cond}}\textsubscript{1} \deleted{Si} Si \deleted{\unsure{p}}parmi deux fig. élémentaires $\deleted{\unsure{Z|F}}X, X'$, on pose
\[
\begin{aligned}
X'\leq X\quad &\text{ssi (déf)}\quad X'\ll X,\ \text{et }X'\text{ et }X\\
&\hphantom{\text{ssi (déf)}\quad X'\ll X,\ }\text{sont comp.}
\end{aligned}
\]
alors on trouve une relation \uline{d'ordre},

% @ Figs 1 and Fig 2
De plus, en termes de cette relation d'ordre, on a les propriétés \uline{Fig} 1

\margin{$\vert\vert\vert$}
$(\Longleftarrow)$ [Si $X', X''\ll X$, \inserted{et si $X,X'$} sont compatibles : $X$, ils sont compatibles entre eux])

\uline{Fig}\textsubscript{2} (dim combinatoire finie) On va
